\hfuzz=10pt
\section*{Abstract}
		Diese Dokumentation stellt die Herleitung der Parameter der T0-Theorie in logischer Reihenfolge vor, mit dem Ziel, Zirkularität auszuschließen. Die Darstellung zeigt, wie die Feinstrukturkonstante $\alpha = 1/137$ aus geometrischen Prinzipien folgen soll, ohne sie vorauszusetzen. Die Herleitungsschritte werden explizit dokumentiert, damit der Vorwurf der Zirkularität Schritt für Schritt geprüft werden kann. Die frühen Kopplungskonstanten-Formeln dieses Dokuments sind als Vorstufen eingeordnet (siehe die Einordnung am Ende des Dokuments).

	\medskip\noindent\textbf{Hinweis zur Fassung.} Die deutsche und die englische Fassung dieses Dokuments weichen in Aufbau und Umfang voneinander ab; bei früheren Überarbeitungen wurden in beiden Fassungen Teile gekürzt, ergänzt oder anders gefasst. Bei Abweichungen gilt die umfangreichere deutsche Fassung; Abschnitte, die nur die englische Fassung enthält, gelten als Ergänzung.


\medskip\noindent\textbf{Statusmarker.} Aussagen mit \textbf{[S]} sind \emph{spekulativ}: ein Hinweis oder eine Vermutung, die im Korpus noch nicht hergeleitet oder numerisch geprüft ist.

	
	\section{Einleitung}
	
	Die T0-Theorie verfolgt den Ansatz, fundamentale physikalische Konstanten nicht als willkürlich zu behandeln, sondern aus der geometrischen Struktur des dreidimensionalen Raums herzuleiten. Die zentrale Behauptung ist, dass die Feinstrukturkonstante $\alpha = 1/137.036$ keine empirische Eingabe darstellt, sondern aus der Raumgeometrie folgt.
	
	Um die Frage der Zirkularität prüfbar zu machen, wird hier die Herleitung der Parameter in logischer Reihenfolge präsentiert, beginnend mit geometrischen Prinzipien und ohne Verwendung experimenteller Werte außer fundamentalen Naturkonstanten; wo an einer Stelle doch ein Messwert eingeht, ist das im Text vermerkt.
\section{Der geometrische Parameter $\xipar$}

\subsection{Herleitung aus fundamentaler Geometrie}

Der universelle geometrische Parameter $\xipar$ setzt sich aus zwei fundamentalen Komponenten zusammen:
\begin{equation}
	\xipar = \frac{4}{3} \times 10^{-4}
\end{equation}

\subsubsection{Die harmonisch-geometrische Komponente: 4/3 als universelle Quarte}

\textbf{4:3 = die Quarte -- ein universelles harmonisches Verhältnis}

Der Faktor 4/3 wird hier als die \textbf{reine Quarte} gedeutet, eines der grundlegenden harmonischen Intervalle:

\begin{equation}
	\frac{4}{3} = \text{Frequenzverhältnis der reinen Quarte}
\end{equation}

Genau wie musikalische Intervalle universal sind:
\begin{itemize}
	\item \textbf{Oktave:} 2:1 (immer, egal ob Saite, Luftsäule, Membran)
	\item \textbf{Quinte:} 3:2 (immer)
	\item \textbf{Quarte:} 4:3 (immer)
\end{itemize}

Diese Verhältnisse sind \textbf{geometrisch/mathematisch}, nicht materialabhängig.

\textbf{Warum ist die Quarte universal?}

Bei einer schwingenden Kugel/Sphäre:
\begin{itemize}
	\item Wenn man sie in 4 gleiche ``Schwingungszonen'' teilt
	\item Verglichen mit 3 Zonen
	\item Ergibt sich das Verhältnis 4:3
\end{itemize}

Das ist \textbf{reine Geometrie}, unabhängig vom Material.

\textbf{Die harmonischen Verhältnisse im Tetraeder:}

Der Tetraeder enthält BEIDE fundamentalen harmonischen Intervalle:
\begin{itemize}
	\item \textbf{6 Kanten : 4 Flächen = 3:2} (die Quinte)
	\item \textbf{4 Ecken : 3 Kanten pro Ecke = 4:3} (die Quarte)
\end{itemize}

\textbf{Die komplementäre Beziehung:}
Quinte und Quarte sind komplementäre Intervalle - zusammen ergeben sie die Oktave:
\begin{equation}
	\frac{3}{2} \times \frac{4}{3} = \frac{12}{6} = 2 \quad \text{(Oktave)}
\end{equation}

Das Dokument liest dies als harmonische Struktur des Raums:
\begin{itemize}
	\item Der Tetraeder enthält beide fundamentalen Intervalle
	\item Die Quarte (4:3) und Quinte (3:2) sind reziprok komplementär
	\item Die harmonische Struktur ist in sich konsistent
\end{itemize}

\textbf{Weitere Erscheinungen der Quarte in der Physik:}
\begin{itemize}
	\item Kristallgittern (4-fach Symmetrie)
	\item Sphärischen Harmonischen
	\item Der Kugelvolumenformel: $V = \frac{4\pi}{3}r^3$
\end{itemize}

\textbf{Deutung:}
\begin{itemize}
	\item \textbf{Parallele zu Pythagoras:} ``Alles ist Zahl und Harmonie''
	\item \textbf{Der Raum selbst} hat eine harmonische Struktur
	\item \textbf{Teilchen} sind ``Töne'' in dieser kosmischen Harmonie
\end{itemize}

Die T0-Theorie deutet den Raum damit als harmonisch strukturiert, mit 4/3 (der Quarte) als Grundsignatur. Diese Deutung motiviert die Setzung des Faktors 4/3; sie ist keine Herleitung im strengen Sinn.

\textbf{Der Faktor $10^{-4}$:}

\textbf{Schritt-für-Schritt QFT-Herleitung:}

\textbf{1. Loop-Suppression:}
\begin{equation}
	\frac{1}{16\pi^3} = 2.01 \times 10^{-3}
\end{equation}

\textbf{2. T0-berechnete Higgs-Parameter:}
\begin{equation}
	(\lambda_h^{\text{(T0)}})^2 \frac{(v^{\text{(T0)}})^2}{(m_h^{\text{(T0)}})^2} = (0.129)^2 \times \frac{(246.2)^2}{(125.1)^2} = 0.0167 \times 3.88 = 0.0647
\end{equation}

\textbf{3. Fehlender Faktor zu $10^{-4}$:}
\begin{equation}
	\frac{10^{-4}}{2.01 \times 10^{-3}} = 0.0498 \approx 0.05
\end{equation}

\textbf{4. Vollständige Berechnung:}
\begin{equation}
	2.01 \times 10^{-3} \times 0.0647 = 1.30 \times 10^{-4}
\end{equation}

\textbf{Was ergibt $10^{-4}$:}
Es ist der T0-berechnete Higgs-Parameter-Faktor $0.0647 \approx 6.5 \times 10^{-2}$, der die Loop-Suppression um Faktor 20 reduziert:

\begin{equation}
	2.01 \times 10^{-3} \times 6.5 \times 10^{-2} = 1.3 \times 10^{-4}
\end{equation}

Der $10^{-4}$-Faktor entsteht aus: \textbf{QFT-Loop-Suppression} ($\sim 10^{-3}$) $\times$ \textbf{T0-Higgs-Sektor-Suppression} ($\sim 10^{-1}$) $=$ $10^{-4}$.
	\section{Der Massenskalierungsexponent $\kappa$}
	
	Aus der fraktalen Dimension folgt direkt:
	
	\begin{equation}
		\kappa = \frac{D_f^{\,\text{eff}}}{2} = \frac{2.973}{2} \approx 1.487
	\end{equation}
	
	Dieser Exponent bestimmt die nicht-lineare Massenskalierung in der T0-Theorie.
	
	\section{Leptonen-Massen aus Quantenzahlen}
	
	Die Massen der Leptonen folgen aus der fundamentalen Massenformel:
	
	\begin{equation}
		m_x = \frac{\hbar c}{\xi^2} \times f(n, l, j)
	\end{equation}

	wobei $f(n, l, j)$ der geometrische Strukturfaktor der Quantenzahlen ist (keine Frequenz und nicht die Grundwindungszahl $f = 1/\xi$):
	
	\begin{align}
		f(n, l, j) = \sqrt{n(n+l)} \times \left[j + \frac{1}{2}\right]^{1/2}
	\end{align}
	
	Für die drei Leptonen ergibt sich:
	
	\begin{itemize}
		\item Elektron $(n=1, l=0, j=1/2)$: $m_e = 0.511$ MeV
		\item Myon $(n=2, l=0, j=1/2)$: $m_\mu = 105.66$ MeV
		\item Tau $(n=3, l=0, j=1/2)$: $m_\tau = 1776.86$ MeV
	\end{itemize}
	
	Nach dem Anspruch dieses Dokuments sind diese Massen keine empirischen Eingaben, sondern folgen aus $\xi$ und den Quantenzahlen (zur Einordnung siehe das Ende des Dokuments).
	
	\section{Die charakteristische Energie $E_0$}
	
	Die charakteristische Energie $E_0$ folgt aus dem geometrischen Mittel der Elektron- und Myonmasse mit fraktaler Korrektur (folgender Abschnitt): $E_0 = \sqrt{m_e m_\mu/K_{\text{frak}}} = 7.398$ MeV.
	
	\section{Alternative Herleitung von $E_0$ aus Massenverhältnissen}
	
	\subsection{Das geometrische Mittel der Lepton-Energien}
	
	Eine alternative Herleitung von $E_0$ ergibt sich direkt aus dem geometrischen Mittel der Elektron- und Myon-Massen:
	
	\begin{equation}
		E_0 = \sqrt{m_e \cdot m_\mu} \cdot c^2
	\end{equation}
	
	Mit den aus Quantenzahlen berechneten Massen:
	\begin{align}
		E_0 &= \sqrt{0.511 \text{ MeV} \times 105.66 \text{ MeV}}\\
		&= \sqrt{54.00 \text{ MeV}^2}\\
		&= 7.35 \text{ MeV}
	\end{align}
	
	\subsection{Vergleich mit dem verwendeten Wert}
	
	Der Wert aus dem geometrischen Mittel (7.35 MeV) stimmt gut mit dem verwendeten Wert (7.398 MeV) überein. Die Differenz beträgt weniger als 1\%:
	
	\begin{equation}
		\Delta = \frac{7.398 - 7.35}{7.35} \times 100\% = 0.65\%
	\end{equation}
	
	\subsection{Physikalische Interpretation}
	
	Dass $E_0$ dem geometrischen Mittel der Lepton-Energien entspricht, wird hier physikalisch so gedeutet:
	
	\begin{itemize}
		\item $E_0$ repräsentiert eine natürliche elektromagnetische Energieskala zwischen Elektron und Myon
		\item Die Beziehung ist rein geometrisch und benötigt keine Kenntnis von $\alpha$
		\item Das Massenverhältnis $m_\mu/m_e = 206.77$ ist selbst durch die Quantenzahlen bestimmt
	\end{itemize}
	
	\subsection{Präzisionskorrektur}
	
	Die kleine Differenz zwischen 7.35 MeV und 7.398 MeV ist die fraktale Korrektur:
	
	\begin{equation}
		E_0^{\text{korrigiert}} = \frac{E_0^{\text{geom}}}{\sqrt{K_{\text{frak}}}} = \frac{7.348}{\sqrt{1 - 100\xi}} = \frac{7.348}{0.99331} = 7.397 \approx 7.398 \text{ MeV}
	\end{equation}
	
	\subsection{Verifikation der Feinstrukturkonstante}
	
	Mit dem geometrisch hergeleiteten $E_0 = 7.35$ MeV:
	
	\begin{align}
		\varepsilon &= \xi \cdot E_0^2\\
		&= (1.333 \times 10^{-4}) \times (7.35)^2\\
		&= (1.333 \times 10^{-4}) \times 54.02\\
		&= 7.20 \times 10^{-3}\\
		&= \frac{1}{138.9}
	\end{align}
	
	Mit dem fraktal korrigierten Wert $E_0 = 7.348/\sqrt{K_{\text{frak}}}$ ergibt sich $\alpha^{-1} = K_{\text{frak}}/(\xi \cdot 7.348^2) = 0.98667 \times 138.91 = 137.06$ ($+0.017\,\%$). Dies stützt die Aussage, dass $E_0$ unabhängig von der Kenntnis der Feinstrukturkonstante hergeleitet werden kann.
	%-----
	
	%-----
	Die beiden Wege zu $E_0$ sind das geometrische Mittel ohne und mit fraktaler Korrektur, $7.348$ und $7.397$ MeV (vorheriger Abschnitt); umgekehrt legt das gemessene $\alpha$ den Faktor $K_{\text{frak}} = 0.98650$ fest, $1.7\times10^{-4}$ neben $1 - 100\xi$ (A130).
	
	\section{Der T0-Kopplungsparameter $\varepsilon$}
	
	Der T0-Kopplungsparameter ergibt sich als:
	
	\begin{equation}
		\varepsilon = \xi \cdot E_0^2
	\end{equation}
	
	Mit den hergeleiteten Werten:
	\begin{align}
		\varepsilon &= (1.333 \times 10^{-4}) \times (7.398 \text{ MeV})^2\\
		&= 7.297 \times 10^{-3}\\
		&= \frac{1}{137.036}
	\end{align}
	
	Die Übereinstimmung mit der Feinstrukturkonstante war nicht vorausgesetzt, sondern ergibt sich als Resultat der geometrischen Herleitung.
	\section*{Die einfachste Formel für die Feinstrukturkonstante}

\[
\boxed{\alpha = \xi \cdot \left(\frac{E_0}{1 \text{ MeV}}\right)^2}
\]
\begin{tcolorbox}[colback=red!5!white,colframe=red!75!black]
	\textbf{Wichtig:} Die Normierung $(1 \text{ MeV})^2$ ist für dimensionslose Ergebnisse notwendig.
\end{tcolorbox}	
	\section{Fraktale Dämpfung: kein unabhängiger zweiter Weg}
	
	Ein scheinbar zweiter Weg zu $\alpha$ führt über eine fraktale Dämpfung:
	
	\begin{equation}
		\alpha_{\text{nackt}}^{-1} = 3\pi \times \xi^{-1} \times \ln\left(\frac{\Lambda_{\text{Planck}}}{m_\mu}\right)
	\end{equation}
	
	Mit dem fraktalen Dämpfungsfaktor:
	\begin{equation}
		D_{\text{frak}} = \left(\frac{\lambda_C^{(\mu)}}{\ell_P}\right)^{D_f-2}
	\end{equation}
	
	soll sich ergeben:
	\begin{equation}
		\alpha^{-1} = \alpha_{\text{nackt}}^{-1} \times D_{\text{frak}}
	\end{equation}
	
	Diese Kette liefert $137.036$ nicht aus eigener Kraft: $\lambda_C^{(\mu)}/\ell_P = 1.156 \times 10^{20}$; hoch $D_f - 2$ ergibt das mit $D_f = 3 - \xi$ rund $1.15 \times 10^{20}$. Mit $\alpha_{\text{nackt}}^{-1} = 3\pi \times 7500 \times \ln(1.2209 \times 10^{22}/105.66) = 3\pi \times 7500 \times 46.196 = 3.265 \times 10^{6}$ müsste $D_{\text{frak}}$ genau $137.036/3.265 \times 10^{6} = 4.2 \times 10^{-5}$ sein; dieser Wert ist an $\alpha$ angepasst (der Exponent müsste $-0.218$ sein). Eine unabhängige Herleitung ist das nicht; tragend ist $\alpha = \xi E_0^2$ mit $E_0 = 7.348/\sqrt{K_{\text{frak}}}$.
	
	\section{Klärung: Die zwei verschiedenen $\kappa$-Parameter}
	
	\subsection{Wichtige Unterscheidung}
	
	In der T0-Theorie-Literatur werden zwei physikalisch unterschiedliche Parameter mit dem Symbol $\kappa$ bezeichnet, was zu Verwirrung führen kann. Diese müssen klar unterschieden werden:
	
	\begin{enumerate}
		\item $\kappa_{\text{mass}} = 1.487$ - Der fraktale Massenskalierungsexponent
		\item $\kappa_{\text{grav}}$ - Der Gravitationsfeldparameter
	\end{enumerate}
	
	\subsection{Der Massenskalierungsexponent $\kappa_{\text{mass}}$}
	
	Dieser Parameter wurde bereits in Abschnitt 4 hergeleitet:
	
	\begin{equation}
		\kappa_{\text{mass}} = \frac{D_f^{\,\text{eff}}}{2} = 1.487
	\end{equation}
	Er ist dimensionslos und bestimmt die Skalierung in der Formel für magnetische Momente:
	
	\begin{equation}
		a_x \propto \left(\frac{m_x}{m_\mu}\right)^{\kappa_{\text{mass}}}
	\end{equation}
	
	\subsection{Der Gravitationsfeldparameter $\kappa_{\text{grav}}$}
	
	Dieser Parameter entsteht aus der Kopplung zwischen dem intrinsischen Zeitfeld und Materie. Die T0-Lagrangedichte lautet:
	
	\begin{equation}
		\mathcal{L}_{\text{intrinsic}} = \frac{1}{2}\partial_\mu T \partial^\mu T - \frac{1}{2}T^2 - \frac{\rho}{T}
	\end{equation}
	
	Die resultierende Feldgleichung:
	
	\begin{equation}
		\nabla^2 T = -\frac{\rho}{T^2}
	\end{equation}
	
	führt zu einem modifizierten Gravitationspotential:
	
	\begin{equation}
		\Phi(r) = -\frac{GM}{r} + \kappa_{\text{grav}} r
	\end{equation}
	
	\subsection{Beziehung zwischen $\kappa_{\text{grav}}$ und fundamentalen Parametern}
	
	In natürlichen Einheiten gilt:
	
	\begin{equation}
		\kappa_{\text{grav}}^{\text{nat}} = \beta_T^{\text{nat}} \cdot \frac{yv}{r_g^2}
	\end{equation}
	
	Mit $\beta_T = 1$ und $r_g = 2Gm_\mu$:
	
	\begin{equation}
		\kappa_{\text{grav}} = \frac{y_\mu \cdot v}{(2Gm_\mu)^2} = \frac{\sqrt{2} m_\mu \cdot v}{v \cdot 4G^2m_\mu^2} = \frac{\sqrt{2}}{4G^2m_\mu}
	\end{equation}
	
	\subsection{Numerischer Wert und physikalische Bedeutung}
	
	In SI-Einheiten:
	
	\begin{equation}
		\kappa_{\text{grav}}^{\text{SI}} \approx 4.8 \times 10^{-11} \text{ m/s}^2
	\end{equation}
	
	Dieser lineare Term im Gravitationspotential:
	\begin{itemize}
		\item Soll die beobachteten flachen Rotationskurven von Galaxien erklären \textbf{[S]}
		\item Würde damit Dunkle Materie entbehrlich machen \textbf{[S]}
		\item Entsteht natürlich aus der Zeitfeld-Materie-Kopplung
	\end{itemize}
	
\section{Vollständige Zuordnung: Standardmodell-Parameter zu T0-Entsprechungen}
\label{sec:sm_t0_mapping}

\subsection{Übersicht der Parameterreduktion}
\label{subsec:parameter_overview}

Das Standardmodell benötigt über 20 freie Parameter, die experimentell bestimmt werden müssen. Das T0-System führt diese auf den einen Parameter $\xi$ zurück, ergänzt durch ganzzahlige Strukturwahlen (rationale Vorfaktoren, Exponenten), die als motivierte Setzungen eingehen, sowie Referenzanker wie $v$ und $E_P$ zur Umrechnung (zur Einordnung der frühen Formeln siehe das Ende des Dokuments):

\begin{equation}
	\boxed{\xi = \frac{4}{3} \times 10^{-4}}
\end{equation}

\subsection{Hierarchisch geordnete Parameter-Zuordnungstabelle}
\label{subsec:hierarchical_mapping}

Die Tabelle ist so organisiert, dass jeder Parameter erst definiert wird, bevor er in nachfolgenden Formeln verwendet wird.

% Tabelle 1: Ebenen 0-3 (Kernparameter)
\begin{table}[htbp]
	\centering
	\adjustbox{max width=\textwidth}{%
\begin{tabular}{p{4cm}p{2.5cm}p{3.5cm}p{2.7cm}}
			\toprule
			\textbf{SM-Parameter} & \textbf{SM-Wert} & \textbf{T0-Formel} & \textbf{T0-Wert} \\
			\midrule
			\multicolumn{4}{l}{\textbf{EBENE 0: FUNDAMENTALE GEOMETRISCHE KONSTANTE}} \\
			\midrule
			Geometrischer Parameter $\xi$ & -- & $\xi = \frac{4}{3} \times 10^{-4}$ & $1.333 \times 10^{-4}$ \\
			& & (aus der Geometrie) & (exakt) \\[0.3em]
			\midrule
			\multicolumn{4}{l}{\textbf{EBENE 1: PRIMÄRE KOPPLUNGSKONSTANTEN (nur von $\xi$ abhängig)}} \\
			\midrule
			Starke Kopplung $\alpha_S$ & $\alpha_S \approx 0.118$ & $\alpha_S = \xi^{-1/3}$ & $19.57$ \\
			& (bei $M_Z$) & $= (1.333 \times 10^{-4})^{-1/3}$ & (nat. Einheiten) \\[0.3em]
			Schwache Kopplung $\alpha_W$ & $\alpha_W \approx 1/30$ & $\alpha_W = \xi^{1/2}$ & $1.15 \times 10^{-2}$ \\
			& & $= (1.333 \times 10^{-4})^{1/2}$ & \\[0.3em]
			Gravitationskopplung $\alpha_G$ & nicht im SM & $\alpha_G = \xi^{2}$ & $1.78 \times 10^{-8}$ \\
			& & $= (1.333 \times 10^{-4})^{2}$ & \\[0.3em]
			Elektromagnetische Kopplung & $\alpha = 1/137.036$ & $\alpha_{EM} = 1$ (Konvention) & $1$ \\
			& & $\varepsilon_T = \xi \cdot E_0^2$ & $7.297 \times 10^{-3}$ \\
			& & (physikalische Kopplung) & (*siehe Anm.) \\[0.3em]
			\midrule
			\multicolumn{4}{l}{\textbf{EBENE 2: ENERGIESKALEN (von $\xi$ und Planck-Skala)}} \\
			\midrule
			Planck-Energie $E_P$ & $1.22 \times 10^{19}$ GeV & Referenzskala & $1.22 \times 10^{19}$ GeV \\
			& & (aus $G, \hbar, c$) & \\[0.3em]
			Higgs-VEV $v$ & $246.22$ GeV & $v = \xi^4 E_P/(5\pi)$ & $245.65$ GeV \\
			& (theoretisch) & (siehe Anhang) & \\[0.3em]
			QCD-Skala $\Lambda_{QCD}$ & $\sim 217$ MeV & $\Lambda_{QCD} = v \cdot \xi^{1/3}$ & $12.6$ GeV \\
			& (freier Parameter) & $= 246 \text{ GeV} \cdot \xi^{1/3}$ & \\[0.3em]
			\midrule
			\multicolumn{4}{l}{\textbf{EBENE 3: HIGGS-SEKTOR (von $v$ abhängig)}} \\
			\midrule
			Higgs-Masse $m_h$ & $125.25$ GeV & $m_h = v \cdot \xi^{1/4}$ & $26.4$ GeV \\
			& (gemessen) & $= 246 \cdot (1.333 \times 10^{-4})^{1/4}$ & \\[0.3em]
			Higgs-Selbstkopplung $\lambda_h$ & $0.13$ & $\lambda_h = \frac{m_h^2}{2v^2}$ & $0.129$ \\
			& (abgeleitet) & $= \frac{(125)^2}{2(246)^2}$ & \\[0.3em]
			\bottomrule
		\end{tabular}%
}
	\caption{Standardmodell-Parameter in hierarchischer Ordnung ihrer T0-Ableitung (Teil 1: Ebenen 0--3)}
	\label{tab:sm-params-1}
\end{table}
Zu Tabelle~\ref{tab:sm-params-1}: Die Formeln $\Lambda_{QCD} = v\cdot\xi^{1/3}$ ($246\times0.05109 = 12.57$~GeV) und $m_h = v\cdot\xi^{1/4}$ ($246\times0.10746 = 26.4$~GeV) treffen die Messwerte ($\sim 217$~MeV bzw.\ $125.25$~GeV) nicht; $\lambda_h$ darunter ist mit dem Messwert $m_h = 125$~GeV gerechnet.
% Tabelle 2a: Ebenen 4-5 (Fermionen-Massen und Neutrinos)
\begin{table}[htbp]
	\centering
	\adjustbox{max width=\textwidth}{%
\begin{tabular}{p{4cm}p{2.5cm}p{3.5cm}p{2.7cm}}
			\toprule
			\textbf{SM-Parameter} & \textbf{SM-Wert} & \textbf{T0-Formel} & \textbf{T0-Wert} \\
			\midrule
			\multicolumn{4}{l}{\textbf{EBENE 4: FERMION-MASSEN (von $v$ und $\xi$ abhängig)}} \\
			\midrule
			\multicolumn{4}{l}{\textit{Leptonen:}} \\
			Elektronmasse $m_e$ & $0.511$ MeV & $m_e = v \cdot \frac{4}{3} \cdot \xi^{3/2}$ & $0.502$ MeV \\
			& (freier Parameter) & $= 246 \text{ GeV} \cdot \frac{4}{3} \cdot \xi^{3/2}$ & \\[0.3em]
			Myonmasse $m_\mu$ & $105.66$ MeV & $m_\mu = v \cdot \frac{16}{5} \cdot \xi^1$ & $105.0$ MeV \\
			& (freier Parameter) & $= 246 \text{ GeV} \cdot \frac{16}{5} \cdot \xi$ & \\[0.3em]
			Taumasse $m_\tau$ & $1776.86$ MeV & $m_\tau = v \cdot \frac{25}{9} \cdot \xi^{2/3}$ & $1783$ MeV \\
			& (freier Parameter) & $= 246 \text{ GeV} \cdot \frac{25}{9} \cdot \xi^{2/3}$ & \\[0.3em]
			\multicolumn{4}{l}{\textit{Up-Typ Quarks:}} \\
			Up-Quarkmasse $m_u$ & $2.16$ MeV & $m_u = v \cdot 6 \cdot \xi^{3/2}$ & $2.27$ MeV \\
			Charm-Quarkmasse $m_c$ & $1.27$ GeV & $m_c = v \cdot 2 \cdot \xi^{2/3}$ & $1.284$ GeV \\
			Top-Quarkmasse $m_t$ & $172.76$ GeV & $m_t = v \cdot \frac{1}{28} \cdot \xi^{-1/3}$ & $173.0$ GeV \\
			\multicolumn{4}{l}{\textit{Down-Typ Quarks:}} \\
			Down-Quarkmasse $m_d$ & $4.67$ MeV & $m_d = v \cdot \frac{25}{2} \cdot \xi^{3/2}$ & $4.72$ MeV \\
			Strange-Quarkmasse $m_s$ & $93.4$ MeV & $m_s = v \cdot 3 \cdot \xi^1$ & $98.4$ MeV \\
			Bottom-Quarkmasse $m_b$ & $4.18$ GeV & $m_b = v \cdot \frac{3}{2} \cdot \xi^{1/2}$ & $4.254$ GeV \\
			\midrule
			\multicolumn{4}{l}{\textbf{EBENE 5: NEUTRINO-MASSEN (von $v$ und doppeltem $\xi$ abhängig)}} \\
			\midrule
			Elektron-Neutrino $m_{\nu_e}$ & $< 2$ eV & $m_{\nu_e} = v \cdot r_{\nu_e} \cdot \xi^{3/2} \cdot \xi^3$ & $\sim 10^{-6}$ eV \\
			& (obere Grenze) & mit $r_{\nu_e} \sim 1$ & (Vorhersage) \\[0.3em]
			Myon-Neutrino $m_{\nu_\mu}$ & $< 0.19$ MeV & $m_{\nu_\mu} = v \cdot r_{\nu_\mu} \cdot \xi^{1} \cdot \xi^3$ & $\sim 10^{-4}$ eV \\
			Tau-Neutrino $m_{\nu_\tau}$ & $< 18.2$ MeV & $m_{\nu_\tau} = v \cdot r_{\nu_\tau} \cdot \xi^{2/3} \cdot \xi^3$ & $\sim 10^{-3}$ eV \\
			\bottomrule
		\end{tabular}%
}
	\caption{Standardmodell-Parameter in hierarchischer Ordnung ihrer T0-Ableitung (Teil 2a: Ebenen 4--5, Fermionen-Massen und Neutrinos)}
	\label{tab:sm-params-2a}
\end{table}
Zur Tabelle oben: Die Koeffizienten $r_\tau = \frac{25}{9}$ und $r_c = 2$ entsprechen Dok.~006 und 046. Für $m_s$ liegt die Formel mit dem Faktor $3$ um $5.4\,\%$ über $93.4$~MeV; Dok.~006 führt $r_s = 26/9$ und damit $94.76$~MeV. Die Neutrinowerte gelten mit $r \sim 1$: $v\,\xi^{3/2}\xi^3 = 9.0\times10^{-7}$~eV, $v\,\xi\,\xi^3 = 7.8\times10^{-5}$~eV, $v\,\xi^{2/3}\xi^3 = 1.5\times10^{-3}$~eV.

% Tabelle 2b: Ebenen 6-7 (Mischungsmatrizen und abgeleitete Parameter)
\begin{table}[htbp]
	\centering
	\adjustbox{max width=\textwidth}{%
\begin{tabular}{p{4cm}p{2.5cm}p{3.5cm}p{2.7cm}}
			\toprule
			\textbf{SM-Parameter} & \textbf{SM-Wert} & \textbf{T0-Formel} & \textbf{T0-Wert} \\
			\midrule
			\multicolumn{4}{l}{\textbf{EBENE 6: MISCHUNGSMATRIZEN (von Massenverhältnissen abhängig)}} \\
			\midrule
			\multicolumn{4}{l}{\textit{CKM-Matrix (Quarks):}} \\
			$|V_{us}|$ (Cabibbo) & $0.22452$ & $|V_{us}| = \sqrt{\frac{m_d}{m_s}} \cdot f_{Cab}$ & $0.21$ \\
			& & mit $f_{Cab} = \sqrt{\frac{m_s - m_d}{m_s + m_d}}$ & \\[0.3em]
			$|V_{ub}|$ & $0.00365$ & $|V_{ub}| = \sqrt{\frac{m_d}{m_b}} \cdot \xi^{1/4}$ & $0.0037$ \\
			$|V_{ud}|$ & $0.97446$ & $|V_{ud}| = \sqrt{1 - |V_{us}|^2 - |V_{ub}|^2}$ & $0.974$ \\
			& & (Unitarität) & \\[0.3em]
			CKM CP-Phase $\delta_{CKM}$ & $1.20$ rad & $\delta_{CKM} = \arcsin(2\sqrt{2}\xi^{1/2}/3)$ & $0.0109$ rad \\
			\multicolumn{4}{l}{\textit{PMNS-Matrix (Neutrinos):}} \\
			$\theta_{12}$ (Solar) & $33.44^\circ$ & $\theta_{12} = \arcsin\sqrt{m_{\nu_1}/m_{\nu_2}}$ & $33.5^\circ$ \\
			$\theta_{23}$ (Atmosphärisch) & $49.2^\circ$ & $\theta_{23} = \arcsin\sqrt{m_{\nu_2}/m_{\nu_3}}$ & $49^\circ$ \\
			$\theta_{13}$ (Reaktor) & $8.57^\circ$ & $\theta_{13} = \arcsin(\xi^{1/3})$ & $2.93^\circ$ \\
			PMNS CP-Phase $\delta_{CP}$ & unbekannt & $\delta_{CP} = \pi(1 - 2\xi)$ & $3.14$ rad \\
			\midrule
			\multicolumn{4}{l}{\textbf{EBENE 7: ABGELEITETE PARAMETER}} \\
			\midrule
			Weinberg-W. $\sin^2\theta_W$ & $0.2312$ & $\tfrac{1{-}\sqrt{1{-}4\alpha_W}}{4}$ & $0.0058$ \\
			& & mit $\alpha_W$ von Ebene 1 & \\[0.3em]
			Starke CP-Phase $\theta_{QCD}$ & $< 10^{-10}$ & $\theta_{QCD} = \xi^{2}$ & $1.78 \times 10^{-8}$ \\
			& (obere Grenze) & & (ausgeschlossen) \\
			\bottomrule
		\end{tabular}%
}
	\caption{Standardmodell-Parameter in hierarchischer Ordnung ihrer T0-Ableitung (Teil 2b: Ebenen 6--7, Mischungsmatrizen und abgeleitete Parameter)}
	\label{tab:sm-params-2b}
\end{table}
	
	Zur Tabelle oben: Die Formel für $|V_{us}|$ ergibt mit den T0-Massen etwa $0.21$ (mit den Massen aus Dok.~006: $0.213$) und trifft den SM-Wert nicht ($-5\,\%$). Für $\delta_{CKM}$ ergibt $\arcsin(2\sqrt2\cdot\sqrt{\xi}/3) = 0.0109$~rad; erst ohne den Faktor $\sqrt{\xi}$ ergäbe sich $\arcsin(2\sqrt2/3) = 1.231$~rad. Für $\theta_{13}$ ergibt $\arcsin(\xi^{1/3}) = 2.93^\circ$; $8.59^\circ$ liefert erst das Argument $(25\xi)^{1/3}$ mit $25\xi = 1/300 = 1/(5|A_5|)$, die Korpusformel aus Dok.~328/340. Die Formel für $\sin^2\theta_W$ ergibt mit $\alpha_W = 0.01155$ den Wert $0.0058$ und trifft den SM-Wert nicht. Die Vorhersage $\theta_{QCD} = \xi^2 = 1.78 \times 10^{-8}$ liegt etwa 178-mal über der in der Tabelle genannten Grenze $10^{-10}$ aus dem elektrischen Dipolmoment des Neutrons und ist damit experimentell ausgeschlossen.

\subsection{Die hierarchische Ableitungsstruktur}
\label{subsec:hierarchical_structure}

Die Tabelle zeigt die klare Hierarchie der Parameterableitung:

\begin{enumerate}
	\item \textbf{Ebene 0}: Nur $\xi$ als fundamentale Konstante
	\item \textbf{Ebene 1}: Kopplungskonstanten direkt aus $\xi$
	\item \textbf{Ebene 2}: Energieskalen aus $\xi$ und Referenzskalen
	\item \textbf{Ebene 3}: Higgs-Parameter aus Energieskalen
	\item \textbf{Ebene 4}: Fermion-Massen aus $v$ und $\xi$
	\item \textbf{Ebene 5}: Neutrino-Massen mit zusätzlicher Unterdrückung
	\item \textbf{Ebene 6}: Mischungsparameter aus Massenverhältnissen
	\item \textbf{Ebene 7}: Weitere abgeleitete Parameter
\end{enumerate}

Jede Ebene verwendet nur Parameter, die in vorherigen Ebenen definiert wurden.

\subsection{Kritische Anmerkungen}
\label{subsec:critical_notes}

\textbf{(*) Anmerkung zur Feinstrukturkonstante:}

Die Feinstrukturkonstante hat im T0-System eine Doppelfunktion:
\begin{itemize}
	\item $\alpha_{EM} = 1$ ist eine \textbf{Einheitenkonvention} (wie $c = 1$)
	\item $\varepsilon_T = \xi \cdot E_0^2 = 7.297 \times 10^{-3}$ ist die \textbf{physikalische EM-Kopplung}
\end{itemize}

\textbf{Einheitensystem:}
Alle T0-Werte gelten in natürlichen Einheiten mit $\hbar = c = 1$. Für experimentelle Vergleiche ist eine Transformation in SI-Einheiten erforderlich.

\section{Kosmologische Parameter: Standardkosmologie ($\Lambda$CDM) vs T0-System}
\label{sec:cosmic_t0_mapping}

\subsection{Grundsätzlicher Unterschied der Ansätze}
\label{subsec:paradigm_shift}

\begin{tcolorbox}[colback=red!5!white,colframe=red!75!black,title=Hinweis: Grundsätzliche Unterschiede]
	Das T0-System postuliert ein \textbf{statisches, ewiges Universum} ohne Urknall, während die Standardkosmologie auf einem \textbf{expandierenden Universum} mit Urknall basiert. Die Parameter sind daher oft nicht direkt vergleichbar, sondern repräsentieren unterschiedliche physikalische Konzepte.

	Der heutige Korpus klammert den kosmischen Sektor bewusst aus, weil auch das Standardmodell ihn nicht herleitet (P39). Die kosmologischen T0-Aussagen dieses Abschnitts sind daher spekulativ \textbf{[S]}.
\end{tcolorbox}

\subsection{Hierarchisch geordnete kosmologische Parameter}
\label{subsec:cosmic_hierarchical_mapping}
% Tabelle 2a: Ebenen 4-5 (Fermionen-Massen und Neutrinos)
\begin{table}[htbp]
	\centering
	\adjustbox{max width=\textwidth}{%
\begin{tabular}{p{4cm}p{2.5cm}p{3.5cm}p{2.7cm}}
			\toprule
			\textbf{SM-Parameter} & \textbf{SM-Wert} & \textbf{T0-Formel} & \textbf{T0-Wert} \\
			\midrule
			\multicolumn{4}{l}{\textbf{EBENE 4: FERMION-MASSEN (von $v$ und $\xi$ abhängig)}} \\
			\midrule
			\multicolumn{4}{l}{\textit{Leptonen:}} \\
			Elektronmasse $m_e$ & $0.511$ MeV & $m_e = v \cdot \frac{4}{3} \cdot \xi^{3/2}$ & $0.502$ MeV \\
			& (freier Parameter) & $= 246 \text{ GeV} \cdot \frac{4}{3} \cdot \xi^{3/2}$ & \\[0.3em]
			Myonmasse $m_\mu$ & $105.66$ MeV & $m_\mu = v \cdot \frac{16}{5} \cdot \xi^1$ & $105.0$ MeV \\
			& (freier Parameter) & $= 246 \text{ GeV} \cdot \frac{16}{5} \cdot \xi$ & \\[0.3em]
			Taumasse $m_\tau$ & $1776.86$ MeV & $m_\tau = v \cdot \frac{25}{9} \cdot \xi^{2/3}$ & $1783$ MeV \\
			& (freier Parameter) & $= 246 \text{ GeV} \cdot \frac{25}{9} \cdot \xi^{2/3}$ & \\[0.3em]
			\multicolumn{4}{l}{\textit{Up-Typ Quarks:}} \\
			Up-Quarkmasse $m_u$ & $2.16$ MeV & $m_u = v \cdot 6 \cdot \xi^{3/2}$ & $2.27$ MeV \\
			Charm-Quarkmasse $m_c$ & $1.27$ GeV & $m_c = v \cdot 2 \cdot \xi^{2/3}$ & $1.284$ GeV \\
			Top-Quarkmasse $m_t$ & $172.76$ GeV & $m_t = v \cdot \frac{1}{28} \cdot \xi^{-1/3}$ & $173.0$ GeV \\
			\multicolumn{4}{l}{\textit{Down-Typ Quarks:}} \\
			Down-Quarkmasse $m_d$ & $4.67$ MeV & $m_d = v \cdot \frac{25}{2} \cdot \xi^{3/2}$ & $4.72$ MeV \\
			Strange-Quarkmasse $m_s$ & $93.4$ MeV & $m_s = v \cdot 3 \cdot \xi^1$ & $98.4$ MeV \\
			Bottom-Quarkmasse $m_b$ & $4.18$ GeV & $m_b = v \cdot \frac{3}{2} \cdot \xi^{1/2}$ & $4.254$ GeV \\
			\midrule
			\multicolumn{4}{l}{\textbf{EBENE 5: NEUTRINO-MASSEN (von $v$ und doppeltem $\xi$ abhängig)}} \\
			\midrule
			Elektron-Neutrino $m_{\nu_e}$ & $< 2$ eV & $m_{\nu_e} = v \cdot r_{\nu_e} \cdot \xi^{3/2} \cdot \xi^3$ & $\sim 10^{-6}$ eV \\
			& (obere Grenze) & mit $r_{\nu_e} \sim 1$ & (Vorhersage) \\[0.3em]
			Myon-Neutrino $m_{\nu_\mu}$ & $< 0.19$ MeV & $m_{\nu_\mu} = v \cdot r_{\nu_\mu} \cdot \xi^{1} \cdot \xi^3$ & $\sim 10^{-4}$ eV \\
			Tau-Neutrino $m_{\nu_\tau}$ & $< 18.2$ MeV & $m_{\nu_\tau} = v \cdot r_{\nu_\tau} \cdot \xi^{2/3} \cdot \xi^3$ & $\sim 10^{-3}$ eV \\
			\bottomrule
		\end{tabular}%
}
	\caption{Standardmodell-Parameter in hierarchischer Ordnung ihrer T0-Ableitung (Teil 2a: Ebenen 4--5)}
	\label{tab:sm-params-2a}
\end{table}
Zur Tabelle oben: Die Koeffizienten $r_\tau = \frac{25}{9}$ und $r_c = 2$ entsprechen Dok.~006 und 046. Für $m_s$ liegt die Formel mit dem Faktor $3$ um $5.4\,\%$ über $93.4$~MeV; Dok.~006 führt $r_s = 26/9$ und damit $94.76$~MeV. Die Neutrinowerte gelten mit $r \sim 1$: $v\,\xi^{3/2}\xi^3 = 9.0\times10^{-7}$~eV, $v\,\xi\,\xi^3 = 7.8\times10^{-5}$~eV, $v\,\xi^{2/3}\xi^3 = 1.5\times10^{-3}$~eV.

% Tabelle 2b: Ebenen 6-7 (Mischungsmatrizen und abgeleitete Parameter)
\begin{table}[htbp]
	\centering
	\adjustbox{max width=\textwidth}{%
\begin{tabular}{p{4cm}p{2.5cm}p{3.5cm}p{2.7cm}}
			\toprule
			\textbf{SM-Parameter} & \textbf{SM-Wert} & \textbf{T0-Formel} & \textbf{T0-Wert} \\
			\midrule
			\multicolumn{4}{l}{\textbf{EBENE 6: MISCHUNGSMATRIZEN (von Massenverhältnissen abhängig)}} \\
			\midrule
			\multicolumn{4}{l}{\textit{CKM-Matrix (Quarks):}} \\
			$|V_{us}|$ (Cabibbo) & $0.22452$ & $|V_{us}| = \sqrt{\frac{m_d}{m_s}} \cdot f_{Cab}$ & $0.21$ \\
			& & mit $f_{Cab} = \sqrt{\frac{m_s - m_d}{m_s + m_d}}$ & \\[0.3em]
			$|V_{ub}|$ & $0.00365$ & $|V_{ub}| = \sqrt{\frac{m_d}{m_b}} \cdot \xi^{1/4}$ & $0.0037$ \\
			$|V_{ud}|$ & $0.97446$ & $|V_{ud}| = \sqrt{1 - |V_{us}|^2 - |V_{ub}|^2}$ & $0.974$ \\
			& & (Unitarität) & \\[0.3em]
			CKM CP-Phase $\delta_{CKM}$ & $1.20$ rad & $\delta_{CKM} = \arcsin(2\sqrt{2}\xi^{1/2}/3)$ & $0.0109$ rad \\
			\multicolumn{4}{l}{\textit{PMNS-Matrix (Neutrinos):}} \\
			$\theta_{12}$ (Solar) & $33.44^\circ$ & $\theta_{12} = \arcsin\sqrt{m_{\nu_1}/m_{\nu_2}}$ & $33.5^\circ$ \\
			$\theta_{23}$ (Atmosphärisch) & $49.2^\circ$ & $\theta_{23} = \arcsin\sqrt{m_{\nu_2}/m_{\nu_3}}$ & $49^\circ$ \\
			$\theta_{13}$ (Reaktor) & $8.57^\circ$ & $\theta_{13} = \arcsin(\xi^{1/3})$ & $2.93^\circ$ \\
			PMNS CP-Phase $\delta_{CP}$ & unbekannt & $\delta_{CP} = \pi(1 - 2\xi)$ & $3.14$ rad \\
			\midrule
			\multicolumn{4}{l}{\textbf{EBENE 7: ABGELEITETE PARAMETER}} \\
			\midrule
			Weinberg-W. $\sin^2\theta_W$ & $0.2312$ & $\tfrac{1{-}\sqrt{1{-}4\alpha_W}}{4}$ & $0.0058$ \\
			& & mit $\alpha_W$ von Ebene 1 & \\[0.3em]
			Starke CP-Phase $\theta_{QCD}$ & $< 10^{-10}$ & $\theta_{QCD} = \xi^{2}$ & $1.78 \times 10^{-8}$ \\
			& (obere Grenze) & & (ausgeschlossen) \\
			\bottomrule
		\end{tabular}%
}
	\caption{Standardmodell-Parameter in hierarchischer Ordnung ihrer T0-Ableitung (Teil 2b: Ebenen 6--7)}
	\label{tab:sm-params-2b}
\end{table}
	
	Zur Tabelle oben: Die Formel für $|V_{us}|$ ergibt mit den T0-Massen etwa $0.21$ (mit den Massen aus Dok.~006: $0.213$) und trifft den SM-Wert nicht ($-5\,\%$). Für $\delta_{CKM}$ ergibt $\arcsin(2\sqrt2\cdot\sqrt{\xi}/3) = 0.0109$~rad; erst ohne den Faktor $\sqrt{\xi}$ ergäbe sich $\arcsin(2\sqrt2/3) = 1.231$~rad. Für $\theta_{13}$ ergibt $\arcsin(\xi^{1/3}) = 2.93^\circ$; $8.59^\circ$ liefert erst das Argument $(25\xi)^{1/3}$ mit $25\xi = 1/300 = 1/(5|A_5|)$, die Korpusformel aus Dok.~328/340. Die Formel für $\sin^2\theta_W$ ergibt mit $\alpha_W = 0.01155$ den Wert $0.0058$ und trifft den SM-Wert nicht. Die Vorhersage $\theta_{QCD} = \xi^2 = 1.78 \times 10^{-8}$ liegt etwa 178-mal über der in der Tabelle genannten Grenze $10^{-10}$ aus dem elektrischen Dipolmoment des Neutrons und ist damit experimentell ausgeschlossen.
\subsection{Kritische Unterschiede und Testmöglichkeiten}
\label{subsec:critical_differences}

\begin{table}[h]
	\centering
	\adjustbox{max width=\textwidth}{%
\begin{tabular}{p{4cm}p{5cm}p{5cm}}
		\toprule
		\textbf{Phänomen} & \textbf{$\Lambda$CDM-Erklärung} & \textbf{T0-Deutung [S]} \\
		\midrule
		Rotverschiebung & Raumexpansion & Fraktale Wegverl\"angerung im $\xi$-Feld (geometrische Wegstreckung; kein Photon-Energieverlust) \\
		CMB & Rekombination bei $z=1100$ ($\Lambda$CDM-Wert) & $\xi$-Feld Gleichgewichtsstrahlung \\
		Dunkle Energie & 68\% des Universums & Nicht erforderlich \\
		Dunkle Materie & 26\% des Universums & $\xi$-Feld Gravitationseffekte \\
		JWST-Paradox & Unerklärte frühe Galaxien & Kein Problem im ewigen Universum \\
		\bottomrule
	\end{tabular}%
}
	\caption{Unterschiede zwischen $\Lambda$CDM und T0 (T0-Spalte spekulativ, vgl.\ P39)}
\end{table}

\subsection{Kritische Anmerkungen zur Testbarkeit}
\label{subsec:testability_notes}

\textbf{(*) Zur Rotverschiebung:}

Die Rotverschiebung ist im T0-System achromatisch (fraktale Wegverlängerung, siehe Tabelle oben); ein Test über eine Wellenlängenabhängigkeit der Rotverschiebung entfällt damit.

\textbf{(**) Zur Masse-Energie-Äquivalenz:}

Die Formel $E = mc^2$ gilt in beiden Systemen identisch. Der Unterschied liegt in der \textbf{Interpretation}:

\begin{itemize}
	\item \textbf{$\Lambda$CDM}: Masse ist eine fundamentale Eigenschaft der Teilchen
	\item \textbf{T0-System}: Masse entsteht durch Resonanzen im $\xi$-Feld (siehe Yukawa-Parameter-Herleitung)
\end{itemize}

Die Formel selbst bleibt unverändert, aber im T0-System ist $m$ keine Konstante, sondern $m = m(\xi, E_{field})$ - eine Funktion der Feldgeometrie. Praktisch macht das keinen messbaren Unterschied für $E = mc^2$.
\appendix

\section{Anhang: Rein theoretische Ableitung des Higgs-VEV aus Quantenzahlen}

\subsection{Fundamentale theoretische Grundlagen}

\subsubsection{Quantenzahlen der Leptonen in der T0-Theorie}

Die T0-Theorie ordnet jedem Teilchen Quantenzahlen $(n, l, j)$ zu, die aus der Lösung der dreidimensionalen Wellengleichung im Energiefeld entstehen:

\textbf{Elektron (1. Generation):}
\begin{itemize}
	\item Hauptquantenzahl: $n = 1$
	\item Bahndrehimpuls: $l = 0$ (s-artig, sphärisch symmetrisch)
	\item Gesamtdrehimpuls: $j = 1/2$ (Fermion)
\end{itemize}

\textbf{Myon (2. Generation):}
\begin{itemize}
	\item Hauptquantenzahl: $n = 2$
	\item Bahndrehimpuls: $l = 1$ (p-artig, Dipolstruktur)
	\item Gesamtdrehimpuls: $j = 1/2$ (Fermion)
\end{itemize}

\subsubsection{Universelle Massenformeln}

Die T0-Theorie verwendet zwei Formulierungen für Teilchenmassen; wie die Massenverhältnisse unten zeigen, sind sie nicht äquivalent:

\textbf{Direkte Methode:}
\begin{equation}
	m_i = \frac{1}{\xi_i} = \frac{1}{\xi_0 \times f(n_i, l_i, j_i)}
	\label{eq:direct_mass_formula}
\end{equation}

\textbf{Erweiterte Yukawa-Methode:}
\begin{equation}
	m_i = y_i \times v
	\label{eq:yukawa_mass_formula}
\end{equation}

wobei:
\begin{itemize}
	\item $\xi_0 = \frac{4}{3} \times 10^{-4}$: Universeller geometrischer Parameter
	\item $f(n_i, l_i, j_i)$: Geometrische Faktoren aus Quantenzahlen
	\item $y_i$: Yukawa-Kopplungen
	\item $v$: Higgs-VEV (Zielgröße)
\end{itemize}

\subsection{Theoretische Berechnung der geometrischen Faktoren}

\subsubsection{Geometrische Faktoren aus Quantenzahlen}

Die geometrischen Faktoren ergeben sich aus der analytischen Lösung der dreidimensionalen Wellengleichung. Für die fundamentalen Leptonen:

\textbf{Elektron $(n=1, l=0, j=1/2)$:}

Die Grundzustandslösung der 3D-Wellengleichung liefert den einfachsten geometrischen Faktor:
\begin{equation}
	f_e(1,0,1/2) = 1
\end{equation}

Dies ist die Referenzkonfiguration (Grundzustand).

\textbf{Myon $(n=2, l=1, j=1/2)$:}

Für die erste angeregte Konfiguration mit Dipolcharakter ergibt die Lösung:
\begin{equation}
	f_\mu(2,1,1/2) = \frac{16}{5}
\end{equation}

Dieser Faktor berücksichtigt:
\begin{itemize}
	\item $n^2 = 4$ (Energieniveau-Skalierung)
	\item $\frac{4}{5}$ (l=1 Dipolkorrektur vs. l=0 sphärisch)
\end{itemize}

\subsubsection{Verifikation der Faktoren}

Die geometrischen Faktoren müssen konsistent mit der universellen T0-Struktur sein:

\begin{align}
	\xi_e &= \xi_0 \times f_e = \frac{4}{3} \times 10^{-4} \times 1 = \frac{4}{3} \times 10^{-4}\\
	\xi_\mu &= \xi_0 \times f_\mu = \frac{4}{3} \times 10^{-4} \times \frac{16}{5} = \frac{64}{15} \times 10^{-4}
\end{align}

\subsection{Ableitung der Massenverhältnisse}

\subsubsection{Theoretisches Elektron-Myon-Massenverhältnis}

Mit den geometrischen Faktoren folgt aus der direkten Methode:

\begin{align}
	\frac{m_\mu}{m_e} &= \frac{\xi_e}{\xi_\mu} = \frac{f_e}{f_\mu} = \frac{1}{\frac{16}{5}} = \frac{5}{16}
\end{align}

\textbf{Achtung:} Dies ist das umgekehrte Verhältnis. Da $\xi \propto 1/m$, erhalten wir:

\begin{align}
	\frac{m_\mu}{m_e} &= \frac{f_\mu}{f_e} = \frac{\frac{16}{5}}{1} = \frac{16}{5} = 3.2
\end{align}
Die direkte Methode liefert also $m_\mu/m_e = 16/5 = 3.2$, die Yukawa-Methode unten $207.8$; die beiden Formulierungen ergeben verschiedene Verhältnisse und sind nicht äquivalent. Außerdem wird das Myon hier mit $l = 1$ geführt, im Abschnitt \enquote{Leptonen-Massen aus Quantenzahlen} mit $l = 0$.

\subsubsection{Korrektur durch Yukawa-Kopplungen}

Die Yukawa-Methode berücksichtigt zusätzliche quantenfeldtheoretische Korrekturen:

\textbf{Elektron:}
\begin{equation}
	y_e = \frac{4}{3} \times \xi^{3/2} = \frac{4}{3} \times \left(\frac{4}{3} \times 10^{-4}\right)^{3/2}
\end{equation}

\textbf{Myon:}
\begin{equation}
	y_\mu = \frac{16}{5} \times \xi^1 = \frac{16}{5} \times \frac{4}{3} \times 10^{-4}
\end{equation}

\subsubsection{Berechnung des korrigierten Verhältnisses}

\begin{align}
	\frac{y_\mu}{y_e} &= \frac{\frac{16}{5} \times \frac{4}{3} \times 10^{-4}}{\frac{4}{3} \times \left(\frac{4}{3} \times 10^{-4}\right)^{3/2}}\\
	&= \frac{\frac{16}{5} \times \frac{4}{3} \times 10^{-4}}{\frac{4}{3} \times \frac{4}{3} \times 10^{-4} \times \sqrt{\frac{4}{3} \times 10^{-4}}}\\
	&= \frac{\frac{16}{5}}{\frac{4}{3} \times \sqrt{\frac{4}{3} \times 10^{-4}}}\\
	&= \frac{\frac{16}{5}}{\frac{4}{3} \times 0.01155}\\
	&= \frac{3.2}{0.0154} = 207.8
\end{align}

Dieses theoretische Verhältnis von $207.8$ liegt sehr nahe am experimentellen Wert von $206.768$.

\subsection{Higgs-VEV}

Der Higgs-VEV folgt im Korpus aus der Strukturformel
\begin{equation}
	\frac{v}{E_P} = \frac{\xi^4}{5\pi}, \qquad v = 245{,}65~\text{GeV} \quad (-0{,}23\,\%~\text{gegen}~246{,}22~\text{GeV}),
\end{equation}
dimensionslos; der Faktor $10$ in $5\pi = (\pi/2)\cdot10$ ist an den Messwert angepasst, mit dem nackten $v$ der Leiter wird er zu $10\,K_{\text{frak}}\approx\pi^2$ (Dok.~149, 387). Das Massenverhältnis $m_\mu/m_e = 207.8$ ($0.5\,\%$) aus dem vorigen Abschnitt enthält $v$ nicht und bleibt davon unberührt.

	\section{Schlussfolgerung}
	
	Die Herleitung zeigt:
	\begin{enumerate}
		\item Die behandelten Parameter werden auf $\xi$ und geometrische Prinzipien zurückgeführt
		\item Die Feinstrukturkonstante $\alpha = 1/137$ wird hergeleitet, nicht vorausgesetzt
		\item Tragend ist $\alpha = \xi E_0^2$; die fraktale Dämpfungskette ist kein unabhängiger zweiter Weg
		\item Für $E_0$ führen das geometrische Mittel ohne und mit fraktaler Korrektur zu $7.348$ bzw.\ $7.397$ MeV
		\item Die Herleitungskette setzt $\alpha$ nicht voraus
		\item Die Unterscheidung zwischen $\kappa_{\text{mass}}$ und $\kappa_{\text{grav}}$
	\end{enumerate}
	
	Die T0-Theorie vertritt damit die These, dass die fundamentalen Konstanten der Natur keine willkürlichen Zahlen sind, sondern aus der geometrischen Struktur des Universums folgen. Zur Einordnung der frühen Formeln dieses Dokuments siehe den letzten Abschnitt.
% ========================================
% DEUTSCHE VERSION
% ========================================

\appendix
\section{Verzeichnis der verwendeten Formelzeichen}
\label{app:symbols_de}

\subsection{Fundamentale Konstanten}
\adjustbox{max width=\textwidth}{%
\begin{tabular}{lll}
	\toprule
	\textbf{Symbol} & \textbf{Bedeutung} & \textbf{Wert/Einheit} \\
	\midrule
	$\xi$ & Geometrischer Parameter & $\frac{4}{3} \times 10^{-4}$ (dimensionslos) \\
	$c$ & Lichtgeschwindigkeit & $2.998 \times 10^8$ m/s \\
	$\hbar$ & Reduzierte Planck-Konstante & $1.055 \times 10^{-34}$ J·s \\
	$G$ & Gravitationskonstante & $6.674 \times 10^{-11}$ m³/(kg·s²) \\
	$k_B$ & Boltzmann-Konstante & $1.381 \times 10^{-23}$ J/K \\
	$e$ & Elementarladung & $1.602 \times 10^{-19}$ C \\
\end{tabular}%
}

\subsection{Kopplungskonstanten}
\adjustbox{max width=\textwidth}{%
\begin{tabular}{lll}
	\toprule
	\textbf{Symbol} & \textbf{Bedeutung} & \textbf{Formel} \\
	\midrule
	$\alpha$ & Feinstrukturkonstante & $1/137.036$ (SI) \\
	$\alpha_{EM}$ & Elektromagnetische Kopplung & $1$ (nat. Einh.) \\
	$\alpha_S$ & Starke Kopplung & $\xi^{-1/3}$ \\
	$\alpha_W$ & Schwache Kopplung & $\xi^{1/2}$ \\
	$\alpha_G$ & Gravitationskopplung & $\xi^{2}$ \\
	$\varepsilon_T$ & T0-Kopplungsparameter & $\xi \cdot E_0^2$ \\
	\bottomrule
\end{tabular}%
}

\subsection{Energieskalen und Massen}
\adjustbox{max width=\textwidth}{%
\begin{tabular}{lll}
	\toprule
	\textbf{Symbol} & \textbf{Bedeutung} & \textbf{Wert/Formel} \\
	\midrule
	$E_P$ & Planck-Energie & $1.22 \times 10^{19}$ GeV \\
	$E_\xi$ & Charakteristische Energie & $1/\xi = 7500$ (nat. Einh.) \\
	$E_0$ & Fundamentale EM-Energie & $7.398$ MeV \\
	$v$ & Higgs-VEV & $246.22$ GeV \\
	$m_h$ & Higgs-Masse & $125.25$ GeV \\
	$\Lambda_{QCD}$ & QCD-Skala & $\sim 200$ MeV \\
	$m_e$ & Elektronmasse & $0.511$ MeV \\
	$m_\mu$ & Myonmasse & $105.66$ MeV \\
	$m_\tau$ & Taumasse & $1776.86$ MeV \\
	$m_u, m_d$ & Up-, Down-Quarkmasse & $2.16$, $4.67$ MeV \\
	$m_c, m_s$ & Charm-, Strange-Quarkmasse & $1.27$ GeV, $93.4$ MeV \\
	$m_t, m_b$ & Top-, Bottom-Quarkmasse & $172.76$ GeV, $4.18$ GeV \\
	$m_{\nu_e}, m_{\nu_\mu}, m_{\nu_\tau}$ & Neutrinomassen & $< 2$ eV, $< 0.19$ MeV, $< 18.2$ MeV \\
	\bottomrule
\end{tabular}%
}

\subsection{Kosmologische Parameter}
\adjustbox{max width=\textwidth}{%
\begin{tabular}{lll}
	\toprule
	\textbf{Symbol} & \textbf{Bedeutung} & \textbf{Wert/Formel} \\
	\midrule
	$T_{CMB}$ & CMB-Temperatur & $2.725$ K \\
	$z$ & Rotverschiebung & dimensionslos \\
	$\Omega_\Lambda$ & Dunkle-Energie-Dichte & $0.6847$ (ΛCDM), $0$ (T0) \\
	$\Omega_{DM}$ & Dunkle-Materie-Dichte & $0.2607$ (ΛCDM), $0$ (T0) \\
	$\Omega_b$ & Baryonendichte & $0.0492$ (ΛCDM), $1$ (T0) \\
	$\Lambda$ & Kosmologische Konstante & $(1.1 \pm 0.02) \times 10^{-52}$ m$^{-2}$ \\
	$\rho_\xi$ & ξ-Feld-Energiedichte & $E_\xi^4$ \\
	$\rho_{CMB}$ & CMB-Energiedichte & $4.64 \times 10^{-31}$ kg/m³ \\
	\bottomrule
\end{tabular}%
}

\subsection{Geometrische und abgeleitete Größen}
\adjustbox{max width=\textwidth}{%
\begin{tabular}{lll}
	\toprule
	\textbf{Symbol} & \textbf{Bedeutung} & \textbf{Wert/Formel} \\
	\midrule
	$D_f^{\text{eff}}$ & Effektive fraktale Dimension (kumulativ über RG-Lauf) & $2{,}973$ \\
	$\kappa_{mass}$ & Massenskalierungsexponent & $D_f^{\text{eff}}/2 = 1.487$ \\
	$\kappa_{grav}$ & Gravitationsfeldparameter & $4.8 \times 10^{-11}$ m/s² \\
	$\lambda_h$ & Higgs-Selbstkopplung & $0.13$ \\
	$\theta_W$ & Weinberg-Winkel & $\sin^2\theta_W = 0.2312$ \\
	$\theta_{QCD}$ & Starke CP-Phase & $< 10^{-10}$ (exp.), $\xi^2$ (T0, ausgeschlossen) \\
	$\ell_P$ & Planck-Länge & $1.616 \times 10^{-35}$ m \\
	$\lambda_C$ & Compton-Wellenlänge & $\hbar/(mc)$ \\
	$r_g$ & Gravitationsradius & $2Gm$ \\
	$L_\xi$ & Charakteristische Länge & $\xi$ (nat. Einh.) \\
	\bottomrule
\end{tabular}%
}

\subsection{Mischungsmatrizen}
\adjustbox{max width=\textwidth}{%
\begin{tabular}{lll}
	\toprule
	\textbf{Symbol} & \textbf{Bedeutung} & \textbf{Typischer Wert} \\
	\midrule
	$V_{ij}$ & CKM-Matrixelemente & siehe Tabelle \\
	$|V_{ud}|$ & CKM ud-Element & $0.97446$ \\
	$|V_{us}|$ & CKM us-Element (Cabibbo) & $0.22452$ \\
	$|V_{ub}|$ & CKM ub-Element & $0.00365$ \\
	$\delta_{CKM}$ & CKM CP-Phase & $1.20$ rad \\
	$\theta_{12}$ & PMNS Solar-Winkel & $33.44°$ \\
	$\theta_{23}$ & PMNS Atmosphärisch & $49.2°$ \\
	$\theta_{13}$ & PMNS Reaktor-Winkel & $8.57°$ \\
	$\delta_{CP}$ & PMNS CP-Phase & unbekannt \\
	\bottomrule
\end{tabular}%
}

\subsection{Sonstige Symbole}
\adjustbox{max width=\textwidth}{%
\begin{tabular}{lll}
	\toprule
	\textbf{Symbol} & \textbf{Bedeutung} & \textbf{Kontext} \\
	\midrule
	$n, l, j$ & Quantenzahlen & Teilchenklassifikation \\
	$r_i$ & Rationale Koeffizienten & Yukawa-Kopplungen \\
	$p_i$ & Generationsexponenten & $3/2, 1, 2/3, ...$ \\
	$f(n,l,j)$ & Geometrische Funktion & Massenformel \\
	$\rho_{tet}$ & Tetraeder-Packungsdichte & $0.68$ \\
	$\gamma$ & Universeller Exponent & $1.01$ \\
	$\nu$ & Kristallsymmetrie-Faktor & $0.63$ \\
	$\beta_T$ & Zeit-Feld-Kopplung & $1$ (nat. Einh.) \\
	$y_i$ & Yukawa-Kopplungen & $r_i \cdot \xi^{p_i}$ \\
	$T(x,t)$ & Zeitfeld & T0-Theorie \\
	$E_{field}$ & Energiefeld & Universelles Feld \\
	\bottomrule
\end{tabular}%
}

% ============================================================
% Einordnung der frühen Kopplungskonstanten-Formeln
% ============================================================

\section*{Einordnung der frühen Kopplungskonstanten-Formeln}

Die früh entwickelten Kopplungskonstanten-Formeln in diesem Dokument
sind als Vorstufen einzuordnen, nicht als kanonische FFGFT-Ableitungen.
Maßgeblich sind die späteren Dokumente~160 (Lebensdauerstruktur),
318 (Geltungsbereich der Masseableitungen) und die Galois-Reihe
(Dok.~336--348). Numerische Verwendung der Formeln dieses Dokuments
erfordert Rückgriff auf diese Nachfolgedokumente.
